\documentclass[12pt]{article}
\usepackage[utf8]{inputenc}
\usepackage[T1]{fontenc}
\usepackage{graphicx} % Required for inserting images
%packages utilisés pour les équations
\usepackage{mathptmx}
\usepackage{mathtools}
\usepackage{mathrsfs}
\usepackage{amsmath}%Pour équations multi-lignes
\usepackage{mathbbol}
\usepackage{envmath}%Pour système d'équations
\usepackage{esint}
\usepackage{amsthm}
\usepackage{appendix}%Pour les annexes
\usepackage{enumerate}%Pour les énumérations
\usepackage{hyperref}%Pour les liens
\usepackage{wrapfig}
\usepackage[paper=a4paper,hmargin=2cm,vmargin=2cm]{geometry}
\usepackage[backend=biber,
bibstyle=apa,citestyle=apa,sorting=anyt]{biblatex}%Pour la biblio C'EST ICI QU'ON MODIFIE LE FORMAT A MODIFIER

\usepackage{csquotes}%pour les citations
\addbibresource{Biblio.bib}%Pour la biblio
\usepackage{xspace}%commandes personnalisees
\usepackage[french,english]{babel}
\usepackage{lipsum} %Pour tester la mise en page
\usepackage{subcaption}%placement d'image
\usepackage{multicol}%
\usepackage{multirow}%POur écrire du texte sur plusieurs lignes dans un tableau
%commandes utiles

\numberwithin{equation}{section}%Pour numéroter les équations par section


\title{ pyTQG documentation }
\author{personne30003}
%\AtBeginDocument{\vspace*{1cm minus 1cm}}
%\AtEndDocument{\vspace{1cm minus 1cm}}
\begin{document}
\selectlanguage{english}
\section{Models}
\verb|BarotropicQG|, \verb|ThermalQG|, and \verb|ThermalQG_coupled| classes are instances of the parent class \verb|QG_model|. Their methods are therefore similar. \verb|beta| ($\beta$-plan approcimation) is the only common attributes.
\subsection{Common functions and attributes}
\begin{itemize}
	\item \verb|psi_from_PV(PV, assign = False)|
	\item \verb|PV_from_psi(psi, assign = False)|
	\item \verb|vort_from_PV(PV=None)|
	\item \verb|PV_from_vort(vort = None, assign = False)|
	\item \verb|psi_from_U(U, constant = 0.0)|
	\item \verb|vort_from_U(U)|
\end{itemize}

\subsection{Barotropic QG (class \texttt|BarotropicQG|)}
\subsubsection{Equations}
\begin{align}
	\left(\frac{\partial^2}{\partial x^2} + \frac{\partial^2}{\partial y^2}\right) \psi - \frac{\psi}{R_d^2} &= q -\beta y \\
	\frac{\partial q}{\partial t} &= -\mathrm{J}(\psi, q)
\end{align}
$psi$ : streamfunction, $q$ : potential vorticity. Fields variables : 
\begin{itemize}
	\item streamfunction : \verb|psi| = $\psi$
	\item potential vorticity : \verb|PV| = $q$ :
	\item relative vorticity : \verb|vorticity| = $\zeta = \left(\frac{\partial^2}{\partial x^2} + \frac{\partial^2}{\partial y^2}\right) \psi = q -\beta y + \frac{\psi}{R_d^2}$
	  
\end{itemize}
\subsubsection{Diagnostics}
\begin{itemize}
	\item Total potential vorticity : \verb|PV_total| = $Q = \iint_S q(x, y, t) \mathrm{d}x\mathrm{d}y$
	\item Kinetic energy : \verb|kinetic_energy| = $E_c = \frac{1}{2} \iint_S \left((\frac{\partial \psi}{\partial x})^2 + (\frac{\partial \psi}{\partial y})^2\right) \mathrm{d}x\mathrm{d}y$
	\item Available potential energy : \verb|potential_energy| = $E_p = \frac{1}{2}\iint_S \frac{\psi^2}{R_d^2}\mathrm{d}x\mathrm{d}y$
	\item Total energy : \verb|total_energy| = $E = E_c + E_p$
	\item Enstrophy : \verb|enstrophy| = $Z =\frac{1}{2} \iint_S q^2 \mathrm{d}x\mathrm{d}y$
\end{itemize}
\subsubsection{Parameters}
\begin{itemize}
	\item Rossby Radius : \verb|inv_Rd2| = $\frac{1}{R_d^2}$
	\item Mean zonal transport (for zonal channel configuration) : \verb|T_0| = $T_0$
\end{itemize}
\subsection{Thermal Quasi-Geostrophic (class \texttt{ThermalQG})}
Formulation based on \cite{Warneford2013}
\subsubsection{Equations}
\begin{align}
	\left(\frac{\partial^2}{\partial x^2} + \frac{\partial^2}{\partial y^2}\right) \psi - \frac{\psi}{R_d^2} &= q -\frac{\theta}{R_d^2} -\beta y \\
	\frac{\partial q}{\partial t} &= -\mathrm{J}(\psi, q) + \frac{1}{R_d^2}\mathrm{J}(\psi, \theta) \\
	\frac{\partial \theta}{\partial t} = -\mathrm{J}(\psi, \theta)
\end{align}
$psi$ : streamfunction, $q$ : potential vorticity, $\theta$ : buoyancy. Fields variables : 
\begin{itemize}
	\item streamfunction : \verb|psi| = $\psi$
	\item buoyancy : \verb|theta| = $\theta$
	\item potential vorticity : \verb|PV| = $q$
	\item relative vorticity : \verb|vorticity| = $\zeta = \left(\frac{\partial^2}{\partial x^2} + \frac{\partial^2}{\partial y^2}\right) \psi = q -\beta y + \frac{\psi - \theta}{R_d^2}$
	
\end{itemize}
\subsubsection{Diagnostics}
\begin{itemize}
	\item Total potential vorticity : \verb|PV_total| = $Q = \iint_S q(x, y, t) \mathrm{d}x\mathrm{d}y$
	\item Kinetic energy : \verb|kinetic_energy| = $E_c = \frac{1}{2} \iint_S \left((\frac{\partial \psi}{\partial x})^2 + (\frac{\partial \psi}{\partial y})^2\right) \mathrm{d}x\mathrm{d}y$
	\item Available potential energy : \verb|potential_energy| = $E_p = \frac{1}{2}\iint_S \frac{\psi^2}{R_d^2}\mathrm{d}x\mathrm{d}y$
	\item Total energy : \verb|total_energy| = $E = E_c + E_p$
	\item Total buoyancy : \verb|theta_total| = $\Theta = \iint_S \theta(x, y, t)\mathrm{d}x\mathrm{d}y$
	\item Enstrophy : \verb|enstrophy| = $Z =\frac{1}{2} \iint_S q^2 \mathrm{d}x\mathrm{d}y$
\end{itemize}

\subsubsection{Parameters}

\begin{itemize}
	\item Rossby Radius : \verb|inv_Rd2| = $\frac{1}{R_d^2}$
	\item Mean zonal transport (for zonal channel configuration) : \verb|T_0| = $T_0$
\end{itemize}

\subsection{Coupled Thermal Quasi-Geostrophic (class \texttt{ThermalQG\_coupled})}
\subsubsection{Equations}
\subsubsection{Diagnostics}
\subsubsection{Parameters}
\section{General parameters (class \texttt{Params})}
\section{Numeric}
\section{Time schemes}

We solve ODE like this
\begin{align}
	\frac{\mathrm{d}u(t)}{\mathrm{d}t} &= \mathrm{RHS}(u, t)\\
	u(0) = u^0
\end{align}
Time Step is noted $\Delta t$.
\paragraph{Euler}~\\
\begin{equation}
	u^{n+1} = u^n+\Delta t\times \mathrm{RHS}(u^n, t^n)
\end{equation}
\paragraph{RK2}~\\
\begin{align}
	u^{n+1} &= u^n+\frac{\Delta t}{2}(k_1+k_2)\\
	k_1 &= \mathrm{RHS}(u^n, t^n)\\
	k_2 &= \mathrm{RHS}(u^n +\Delta t k_1, t^n+\Delta t)
\end{align}
\paragraph{RK4}~\\
\begin{align}
	u^{n+1} &= u^n+\frac{\Delta t}{3}(\frac{k_1}{2}+k_2+k_3+\frac{k_4}{2})\\
	k_1 &= \mathrm{RHS}(u^n, t^n)\\
	k_2 &= \mathrm{RHS}(u^n +\frac{\Delta t}{2}k_1, t^n+\frac{\Delta t}{2})\\
	k_3 &= \mathrm{RHS}(u^n+\frac{\Delta t}{2}k_2, t^n+\frac{\Delta t}{2})\\
	k_4 &= \mathrm{RHS}(u^n+\Delta t k_3, t^n+\Delta t)
\end{align}

\paragraph{Leap-Frog}~\\
\begin{align}
	u^{n+1} = u^{n-1}+2\Delta t \mathrm{RHS}(u^n, t^n)
\end{align}
First step is computed with Euler scheme. To remove computationnal mode, we use Asselin filter : 
\begin{align}
	u^{n+1} &= u^{n-1}+2\Delta t \mathrm{RHS}(\overline{u^n}, t^n)\\
	\overline{u^n} &= u^{n} + \alpha(u^{n+1}+2u^n+u^{n-1})\\
	\alpha &\approx 0.05
\end{align}

\printbibliography
\end{document}